\chapter{Integration der Webframeworks in AEM}
\label{sec:integration-der-webframeworks-in-aem}
In diesem Kapitel werden verschiedene Lösungsansätze beschrieben, ein Webframework in eine mit \ac{aem} erstellte Webpräsenz zu integrieren. Dies umfasst auch die Aufbereitung der Webanwendung vor der Integration und Lösungen für Probleme, welche nur indirekt mit der Integration zusammenhängen. Zum Ende dieses Kapitels erfolgt ein Vergleich der Lösungen mit einem kurzen Ausblick darauf, wie diese sich auch miteinander kombinieren ließen.


\section{Rahmenbedingungen}
Der Versuch der Integration erfolgt in ein \ac{aem} der Version 6.0.0.20140515, ausgeführt unter Windows 10 x64 Version 1607, Build 14393.479 mit der \ac{jre} Version 1.8.0\_111-b14.

\section{Zu untersuchende Frameworks}

Im Folgenden soll der Fokus auf die vier Webframeworks aus \autoref{sec:auswahl-und-bewertung} liegen. Diese sind AngularJS (\autoref{sec:angularjs}),  AngularJS2 (\autoref{sec:angularjs2}), Aurelia (\autoref{sec:aurelia}) und React (\autoref{sec:react}).
%Grund für diese Auswahl ist eine zuvor festgelegte Rahmenbedingung innerhalb dieser Arbeit.

\section{Zu integrierende Webanwendungen}
\label{sec:zu-integrierende-webanwendungen}

Als Basis für den Versuch der Integration wurde für jedes Webframework eine Webanwendung ausgewählt. Diese stammen größtenteils direkt von den Entwicklern des jeweiligen Webframeworks und gebrauchen bereits zahlreiche der jeweils verfügbaren Funktionalitäten. Der Quellcode liegt als eigens eingerichtetes GIT Respository zur Einsicht vor. Eine Übersicht der Webframeworks bietet \autoref{tab:webanwendungen}.

\begin{minipage}{\textwidth}
\begin{longtable}{| c | c | c | c |} 
	\hline
	\thead{Framework} & \thead{Basiert auf} & \thead{GIT Respository} \\ 
	
	\hhline{|=|=|=|=|} 
	AngularJS 1.5.8 & angular-phonecat \cite{Angular2016} &  \cite{Kandler2016a} \\
	\hline
	AngularJS2 2.0.0-rc.3& Tutorial: Tour of Heroes \cite{Google2016e} & \cite{Kandler2016b}\\ 
	\hline
	Aurelia 1.1.0& Quick Start \cite{Eisenberg2017} & \cite{Kandler2016c}\\ 
	\hline
	React 15.2.0 & Redux Tetris \cite{Lugo2016} & \cite{Kandler2016d}\\ 
	
	\hline 
	\caption{Webanwendungen}\label{tab:webanwendungen}
\end{longtable}
\end{minipage}

Im Folgenden wird vermehrt der Begriff \quotes{die Webanwendung} verwendet. Hiermit ist eine beliebige der vier genannten Webanwendungen genannt. Die Aussage \quotes{Es wird versucht eine Webanwendung in das \ac{aem} zu integrieren} bedeutet somit, dass nacheinander versucht wird die vier in \ref{tab:webanwendungen} genannten Webanwendungen in ein \ac{aem} zu integrieren.
%Auch wenn hier der Singular verwendet wird, sind damit alle vier Webanwendungen gemeint. 



\section{Möglichkeiten der Integration}
\label{sec:integrationen}
Abhängig von verschiedenen Faktoren wurden unterschiedliche Möglichkeiten erarbeitet, um eine Webanwendung in eine \ac{aem}-Instanz als \ajc zu integrieren. Neben dem zugrunde liegenden clientseitigen Webframework ist der wichtigste Entscheidungsfaktor dafür, welche Lösung genutzt werden kann, die IT-Landschaft der Zielplattform. Je nach dessen Konfiguration bieten sich gewisse Lösungen mehr an oder sind auch als mögliche Lösung von vornherein ausgeschlossen. Auch gewisse wünsche des Kunden, wie die Bereitstellung der App-Ressourcen auf einen seperaten Webserver wird berücksichtigt werden.
%Insbesondere das Thema Zugriffsrechte spielt hier eine immense Rolle. Diese könnten so geregelt sein, dass ein zu tiefer Eingriff in das \ac{aem} und eine freie Konfiguration des Servers nicht gestattet sind.

%\newpage\missing{Stark überarbeiten}\definecolor{lightred}{RGB}{255,230,230}\pagecolor{lightred}

\subsectiontbd{Direkte Integration in das AEM}
\label{sec:direkte-integration-in-das-aem}
Sofern der Zugriff auf das Ziel-\ac{aem}, insbesondere auf dessen \ac{jcr}-Strukur, besteht, lässt sich die Webanwendung direkt in selbiges durch die Verwendung von Clientlibs integrieren.\\
Im ersten Schritt wird das gewünschte JavaScript-Framework als Clientlib, wie in \autoref{sec:clientlib} beschrieben, bereitgestellt. Sollte das Framework mehrmals Verwendung finden empfiehlt es sich dies unter /etc/clientlibs zu platzieren. \\
Um die Bereitstellung der App-Ressourcen für die Webanwendung zu ermöglichen existieren mehrere Ansätze. Je nach verwendetem Framework lassen sich bestimmte Ansätze leichter erreichen.

\subsubsectiontbd{Verwendung als Clientlib}

Besteht die Webanwendung lediglich aus JavaScript und Stylesheets, bestellt die Möglichkeit diese komplett als Clientlib umzusetzen.
Hierfür wird wie in \autoref{sec:clientlib} beschrieben im \ac{jcr} eine entsprechende Struktur benötigt. Das Übertragen der App-Ressourcen in das \ac{jcr} kann über mehrere Wege erfolgen. \\

\paragraphtbd{Transferieren der App-Ressourcen in das AEM}
Zunächst lässt sich über die \ac{webdav}-Schnittstelle mit einem entsprechenden Programm die Struktur des \ac{jcr} anzeigen und bearbeiten. Neben dem Löschen und Verschieben von Knoten ist hier auch das Hochladen von lokalen Dateien möglich. \\
Sofern die \ac{webdav}-Schnittstelle auf der Zielplattform nicht zugänglich sein sollte, weil diese zum Beispiel deaktiviert wurde, gibt es die Alternative die App-Ressourcen als \ac{crx}-Paket bereit zu stellen. Hierfür werden zunächst wieder alle Inhalte in das \ac{jcr}, zum Beispiel über \ac{webdav}, geladen, dieses mal jedoch in dem vom Entwickler. Anschließend kann über die Autoren-Bedienoberfläche, siehe \autoref{sec:autor_ui}, ein \ac{crx}-Paket erstellt werden. Dieses ist ein Zip-Archiv mit Metainformationen, wie Name und Version des \ac{crx}-Pakets.\\
Das Erstellen eines \ac{crx}-Paketes erfolgt über eine entsprechende Seite des CRXDE Lite. Hier werden besagte Metainformationen gesetzt und alle benötigten Knoten für das \ac{crx}-Paket angegeben. Nun erfolgt der Export des \ac{crx}-Pakets. Ist dies geschehen wird das \ac{crx}-Paket in die Zielplattform importiert. Der Inhalt des \ac{crx}-Paketes wird nun in das \ac{jcr} geschrieben. Bereits bestehende Knoten werden bei diesem Vorgang ggf. überschrieben.

\subsubsectiontbd{Content Ordner}
Jeglicher Inhalt, der unter /content abgelegt wird, ist direkt über einen HTTP-Request aufrufbar. Unter der Standardkonfiguration von \ac{aem} ist die Datei \hl{/content/myapp/index.html} zu erreichen über \hl{http://<Server-Domain>:<Port>/content/myapp/index.app}. 
\subsubsectiontbd{Hybride}
Innerhalb einer Clientlib lassen sich lediglich JavaScript und CSS hinterlegen. Falls eine Webanwendung aber noch weitere App-Ressourcen wie Bilder und HTML beinhaltet, so werden diese unter einen anderen Knoten abgelegt. Es bietet sich an diese unter \hl{/content} abzulegen, da wie zuvor beschrieben der Zugriff auf hier abgelegte Dateien direkt zur Verfügung steht.

%\newpage\pagecolor{white}

\subsection{Integration über einen zusätzlichen Webserver}

Neben der Lagerung der App-Ressourcen in dem \ac{jcr} lässt sich diese ebenfalls auf einen zusätzlichen Webserver ablegen und über eine \ac{aem}-Komponente in eine Webseite des \ac{aem} integrieren.

\subsubsection{Lösungsansatz \quotes{Java-Servlet}}
\label{sec:ansatz-java}

Hier ist die Grundidee, dass sich die App-Ressourcen nicht im \ac{jcr} des \ac{aem} befinden, sondern auf einem separaten Webserver vorliegen. Ziel ist es, dass die Webanwendung unabhängig auf dem separaten Webserver zum Testen lauffähig ist. Über eine \ac{aem}-Komponente wird diese in das \ac{aem} importiert. Die \ac{aem}-Komponente bezieht die App-Ressourcen vom Webserver und passt diese ggf. an, so dass die Webanwendung korrekt in der Webseite des \ac{aem} dargestellt wird, wie in \autoref{img:import} illustriert wird. Da auch diese \ac{aem}-Komponente eine Webanwendung abstrahiert, wird diese im Folgenden ebenfals als \ajc bezeichnet.

\begin{figure}[H]
	\begin{center}
		\includesvg[width=1.0\textwidth]{import}
		\caption{Integration vom Webserver in den AEM Server mittels einer \ajc}
		\label{img:import}
	\end{center}
\end{figure}
Dabei werden die App-Ressourcen zur Laufzeit der \ajc, also wenn die Webseite geladen wird, jeweils erneut vom Webserver geladen. Ein Ablegen in das \ac{jcr} ist nicht vorgesehen. Ein temporäres Zwischenspeichern (Cachen) wäre von Grund auf denkbar ist, wird innerhalb dieser Arbeit jedoch nicht weiter behandelt. \\
Im Folgenden wird der separate Webserver, auf dem sich die App-Ressourcen für die Webanwendung befinden, kurz \quotes{Webserver} und der Server mit der lauffähigen \ac{aem}-Instanz kurz \quotes{\ac{aem}-Server} genannt. Die Integration erfolgt somit vom Webserver in den \ac{aem}-Server. Die Webanwendung ist über \serverB aufrufbar und soll später über den \ac{aem}-Server über \serverA erreichbar sein.\\

\input{inhalt/integration-der-webframeworks-in-aem/moeglichkeiten-der-integration/integration-über-einen-zusätzlichen-webserver/ansatz-java/Erklaerung}
\input{inhalt/integration-der-webframeworks-in-aem/moeglichkeiten-der-integration/integration-über-einen-zusätzlichen-webserver/ansatz-java/Anpassen-des-Webservers}
\input{inhalt/integration-der-webframeworks-in-aem/moeglichkeiten-der-integration/integration-über-einen-zusätzlichen-webserver/ansatz-java/Anpassen-der-Webanwendung}
\input{inhalt/integration-der-webframeworks-in-aem/moeglichkeiten-der-integration/integration-über-einen-zusätzlichen-webserver/ansatz-java/Ablauf-der-AEM-Komponente}

\input{inhalt/integration-der-webframeworks-in-aem/moeglichkeiten-der-integration/integration-über-einen-zusätzlichen-webserver/ansatz-java/Bewertung}

\subsubsection{Lösungsansatz \quotes{JavaScript}}
\label{sec:ansatz-javascript}
Dieser Lösungsansatz verzichtet auf die beiden Java-Ressourcen vom Lösungsansatz "Java-Servlet" und erfüllt das gewünschte Ziel lediglich mit JavaScript, welche sich ebenfalls in einer \ac{aem}-Komponente befindet. Im Gegensatz zur zuvor vorgestellten Variante wird nun der erste Schritt von \autoref{img:java} durch eine JavaScript-Komponente ersetzt. Das Prinzip ist hierbei das gleiche, jedoch gibt es hier die Besonderheit, dass es die zu ladenden App-Ressourcen zu synchronisieren gilt.
%Und zwar ist darauf zu achten, dass die in der HTML-Seite referenzierten App-Ressourcen in der korrekten Reihenfolge geladen werden. \\
%Der Ablauf beginnt wieder damit, dass zunächst die reine HTML-Seite geladen wird. Auch diese wird wieder nach JavaScript-Ressourcen durchsucht, die es zu laden gibt. 

\input{inhalt/integration-der-webframeworks-in-aem/moeglichkeiten-der-integration/integration-über-einen-zusätzlichen-webserver/ansatz-javascript/Synchronisation}
\input{inhalt/integration-der-webframeworks-in-aem/moeglichkeiten-der-integration/integration-über-einen-zusätzlichen-webserver/ansatz-javascript/Bewertung}

\subsection{Lösungsansatz \quotes{Proxy}}
\label{sec:sol_proxy}
Ein weiterer Lösungsansatz wird durch Hinzunahme eines Proxy realisierbar. Durch den Einsatz eines Proxy lassen sich die vom Client angeforderten App-Ressourcen während des HTTP-Request-Response-Zykluses manipulieren.\\
Im allgemeinen ist ein Proxy\index{Proxy} ein Rechner, der als Mittelsmann zwischen Client und Server dient. Der Client stellt hierbei die Anfragen an den Proxy. Der Proxy leitet anschließend die Anfrage weiter an den Server. Client und Server kommunizieren somit nie direkt miteinander, sondern nur indirekt über den Proxy.\\
Im Rahmen dieser Arbeit handelt es sich bei einem Proxy um einen weiteren speziell konfigurierten Webserver. Die Web-Ressourcen werden dabei von dem eigentlichen Webserver über den Proxy an den Client ausgeliefert. Die angefragten App-Ressourcen lassen sich hierbei über den Proxy manipulieren. 

\subsubsection{Erklärung}
Die Lösung sieht hierbei vor, dass ein weiterer Webserver als Proxy unter dem Hostnamen \hostnameP~ so angepasst wird, dass dieser alle Anfragen wie in \autoref{img:proxy} über ein zentrales Skript umleitet. Besagtes Skript kann in einer beliebigen, vom Proxy unterstützten Programmiersprache realisiert werden, wie hier in \acs{php}. Die Konfiguration ließe sich beim Apache HTTP-Server mit dem Modul \quotes{mod\_rewrite} realisieren \cite{Foundation2017a}. Die entsprechende Konfiguration ist \autoref{lst:rewrite} zu entnehmen. \\
\begin{minipage}{\textwidth}
\begin{lstlisting}[style=jcr, caption=Konfiguration des mod\_rewrite Modules für den Proxy-Lösungsansatz, label=lst:rewrite]
RewriteEngine on
RewriteRule ^/app/(.*)$ /proxy.php?url=$1 [L]
\end{lstlisting}
\end{minipage}
Der Ablauf könnte hierbei wie in \autoref{img:proxy} erfolgen. \\
Die Konfiguration definiert, dass, sofern eine App-Ressource angefragt wird, deren Pfad mit \filefolder{/app/} beginnt, die Anfrage über \filefolder{proxy.php} umgeleitet wird. Der restliche Pfad, hier \filefolder{templates/news.html}, wird dem Skript als Parameter mit übergeben.\\
Dieses ermittelt nun die benötigte App-Ressource vom eigentlichen Webserver, der unter dem Hostnamen und Port \pseudourl{www.exmaple.com:3000} zu erreichen ist. Etwaige Abweichungen vom angeforderten Pfad und der Pfad auf dem Webserver lassen sich über das Skript wie im Beispiel bewerkstelligen und korrigieren.\\
Anschließend bearbeitet das Skript die App-Ressource entsprechend und liefert die bearbeitete App-Ressource als Ergebnis zurück. Diese Aktionen bleiben dem Client verborgen. Aus dessen Sicht wird ihm lediglich die angeforderte App-Ressource geliefert. \\

\begin{figure}[H]
	\begin{center}
		{\footnotesize\includesvg[width=1.0\textwidth]{proxy}}
		\caption{Funktionsweise des Proxy}
		\label{img:proxy}
	\end{center}
\end{figure}

Die Lösung ähnelt dem Servlet aus \autoref{sec:ansatz-java}, da auch hier zur Laufzeit Änderungen vorgenommen werden können.\\
Seitens \ac{aem} kommt auch hier eine \ac{aem}-Komponente zum Einsatz, welche die initiale Anfrage an die HTML-Seite der Webanwendung stellt, was in etwa dem ersten Schritt aus \autoref{img:java} entspräche. Die Logik der \ajc vom Lösungsansatz \quotes{Java-Servlet} wird hier auf das Proxy-Skript ausgelagert. Innerhalb der \ac{aem}-Komponente wird lediglich die Start-URL des Proxy definiert.
\subsubsection{Bewertung}
Die Lösung mit einem Proxy ist wie folgt zu bewerten.

\begin{minipage}[t]{0.5\textwidth}
	\textbf{Vorteile:}
	\begin{itemize}
		\item Die Struktur der Webanwendung bleibt erhalten.
		\item Die Webanwendung ist unabhängig von \ac{aem} testbar.
		\item Die entwickelte \ajc lässt sich wiederverwenden.
		\item Die Logik des Skriptes ist weitestgehend frei wählbar und nicht von \ac{aem} abhängig.
	\end{itemize}
\end{minipage}
\begin{minipage}[t]{0.5\textwidth}
	\textbf{Nachteile:}
	\begin{itemize}
		\item Betreiben von zwei bzw. drei Diensten nötig (AEM und Webserver und Proxy).
		\item Ggf. zusätzliche Konfiguration von Webserver und Webanwendung nötig.
		\item Zusätzliche Laufzeit (gering, ggf. vernachlässigbar).
	\end{itemize}
\end{minipage}

\section{Möglichkeiten der Kombination}

Jeder der genannten Lösungen ist generell für sich alleine stehend funktionell und erfüllt richtig angewendet das gewünschte Ziel. Es macht aber auch durchaus Sinn einige Lösungsansätze zu kombinieren.

\subsection{Clientlib und Content-Ordner}
Diese Kombination sieht vor, dass die App-Ressourcen in zwei Gruppen unterteilt werden. Zum einen wären hier alle App-Ressourcen, die bereits in der HTML-Seite referenziert sind, also zu Beginn geladen werden. Diese CSS- und JavaScript-Ressourcen werden als Clientlib realisiert. Zum anderen werden alle weiteren Ressourcen unter \pseudourl{/content} abgelegt. 
Folgendes Beispiel zur Erklärung. Es wird angenommen, dass sich die \ajc unter \pseudourl{http://aem.example.com/content/app} befindet. Über die \ajc wurde bereits die Clientlib und somit die im ersten Schritt benötigten App-Ressourcen geladen. Nun wird über Ajax versucht eine JavaScript-Ressource unter der relativen \ac{url} \pseudourl{./templates/app.html} aufzurufen. Somit ergibt sich die absolute Adresse \pseudourl{http://aem.example.com/content/app/templates/app.html}, was innerhalb des \ac{jcr} die Knoten \pseudourl{/content/app/templates/app.html}. Der genaue Ablageknoten unter \pseudourl{/content} ist somit von der \ac{url} abhängig, unter dem die Webanwendung auffindbar sein soll.

\subsubsection{Bewertung}

Die Lösungskombination ist wie folgt zu bewerten.

\begin{minipage}[t]{0.5\textwidth}
	\textbf{Vorteile:}
	\begin{itemize}
		\item Schnell zu realisieren.
		\item Keinerlei Anpassung der Webanwendung nötig. Relative Pfade werden korrekt aufgelöst.
		\item Die Realisierung erfolgt als AJC, somit lässt sich diese frei innerhalb einer Webseite platzieren.
	\end{itemize}
\end{minipage}
\begin{minipage}[t]{0.5\textwidth}
	\textbf{Nachteile:}
	\begin{itemize}
		\item App-Ressoucen müssen innerhalb eines entsprechenden Knoten unter \pseudourl{/content} liegen. Wird die Anwendung in eine andere Webseite eingebettet, müssen die App-Ressourcen in einen anderen Knoten verschoben werden.
	\end{itemize}
\end{minipage}
\subsection{Proxy und Java}

Der Lösungsansatz \quotes{Java-Servlet} hat den Nachteil, dass für jedes Webframework, oder gar für jede Webanwendung eine eigene \ajc benötigt wird. Denn je nach Art der Webanwendung muss die \ajc diese unterschiedlich vor der Integration bearbeiten.  \\
Weiterhin vom Nachteil ist, dass die Webwendung über den Proxy für Besucher voll zugänglich ist. \\
Um die jeweiligen Vorteile zu vereinen werden beide Lösungsansätze hier vereint. 
\begin{figure}[H]
	\begin{center}
		\includesvg[width=.8\textwidth]{proxy+java}
		\caption{Aufbau der Proxy und Java Kombination}
		\label{img:proxyjava}
	\end{center}
\end{figure}

Hierbei wird die Logik, welche für die Anpassung der Webanwendung benötigt wird, in den Proxy verschoben. Somit genügt eine \ac{aem}-Komponente, in welcher nur die Einstiegsseite der Webanwendung zu konfigurieren ist. Die Dateien \filefolder{AppBuilder.java}, \filefolder{GetData.java} und \filefolder{ajaxredirct.js} werden auch hier eingesetzt, dienen jedoch nur zum Auslösen der HTTP-Anfragen.\\
Für jede Webanwendung lässt sich nun auf dem Server ein eigenes Proxy-Skript erstellen, welches sich z. B. bei Apache durch \quotes{mod\_rewrite} einem entsprechenden \ac{url}-Muster zuordnen lässt. \\
Neben der Modularität und das nur eine \ac{aem}-Komponente Verwendung findet besteht die Kommunikation hier nur zwischen der \ajc und dem Webserver. Der Client kommuniziert nie direkt mit dem Proxy. Somit lässt sich der Webserver konfigurieren, dass nur der \ac{aem}-Server und ggf. einige Test-Nutzer auf diesen Zugriff haben. Allen anderen Nutzern wird der Zugriff verweigert.

\subsubsection{Bewertung}

Die Lösungskombination ist wie folgt zu bewerten.

\begin{minipage}[t]{0.5\textwidth}
	\textbf{Vorteile:}
	\begin{itemize}
		\item Nur eine \ac{aem}-Komponente.
		\item Striktere Trennung. Alle, was die Webanwendung betrifft, auch die Logik der Integration, ist außerhalb von \ac{aem}.
		\item Jeder Proxy kann in einer eigenen Skriptsprache erfolgen.
		\item Ggf. weitere Logik in der AEM-Komponente möglich, wie z. B. eine zentrale Komprimierung der vom Proxy erhaltenen Ressourcen.
	\end{itemize}
\end{minipage}
\begin{minipage}[t]{0.5\textwidth}
	\textbf{Nachteile:}
	\begin{itemize}
		\item Externer Webserver wird benötigt.
	\end{itemize}
\end{minipage}

\section{Weitere Werkzeuge und Vorgehensprozesse}
\label{sec:aufbereiten}
Neben der Integration gibt es noch Aufgaben, die es vor und nach der Integration zu erfüllen gilt. So besteht vor der Integration die Möglichkeit, die entwickelte Webanwendung entsprechend aufzubereiten. Dies kann die Performance verbessern und den Prozess der Integration vereinfachen. In diesem Schritt finden Werkzeuge wie Grunt und Gulp aus \autoref{sec:automatiserungswerkzeuge} ihre Verwendung. \\
Weiterhin sollte gerade bei einer \ac{spa} der \ac{aem}-Server entsprechend konfiguriert werden, um Suchmaschinen die korrekte Indizierung der einzelnen Webseiten zu gestatten.

\subsection{Minimierung}
Für den Produktiveinsatz lassen sich die JavaScript-Ressourcen minimieren. Hierbei werden zum Beispiel unnötige Leerzeichen entfernt oder auch lange Variablennamen durch kürzere ersetzt. Dies minimiert die Übertragungszeit vom Server zum Client, da das zu übertragende Volumen sinkt.\\
Generell existieren für die Minimierung von JavaScript-Quellcode verschiedenste Werkzeuge. Eines davon wäre UglifyJS \cite{Bazon2016}. Hiermit lässt sich JavaScript-Quellcode unter anderem minimieren. Das Werkzeug wird unter der BSD-Lizenz veröffentlicht und über \ac{npm} installiert.	
\subsection{Konkatinierung und weitere Laufzeitoptimierung}
\label{sec:konkat}
Für den Produktiveinsatz lassen sich die Quelldateien konkatenieren. Hierbei wird die Anzahl der einzelnen Dateien minimiert, indem diese zu einer oder mehreren großen Dateien zusammengefügt werden. Dies minimiert die Übertragungszeit vom Server zum Client, da weniger HTTP-Request-Response-Zyklen durchlaufen werden müssen. Weiterhin existieren so auch weniger Dateien, die es in das \ac{aem} zu integrieren gilt, was zum Beispiel die direkte Integration, wie in \autoref{sec:direkte-integration-in-das-aem} beschrieben, erleichtert. \\
Weiterhin lassen sich bei manchen Frameworks die HTML-Templates in JavaScript-Dateien umwandeln. Somit ist auch hier eine anschließende Minimierung und Komprimierung möglich.


\begin{description}
	\item[Allgemein] Generell existieren für die Konkatinierung von JavaScript-Quellcodes verschiedenste Werkzeuge. Auch hierfür lässt sich UglifyJS verwenden. Bei der Konkatenierung ist bei dessen Einsatz darauf zu achten, die JavaScript-Dateien in der korrekten Reihenfolge anzugeben, um Fehler in der Abhängigkeit zu vermeiden. \\
	Um solche Abhängigkeitsfehler zu vermeiden lässt sich hier das Konzept der \ac{di} nutzen. Wurde die Webanwendung unter Verwendung von \ac{di} entwickelt, so ist in den JavaScript-Dateien bereits deren Abhängigkeit syntaktisch hinterlegt. Je nach Art der verwendeten \ac{di} darf der Entwickler auf unterschiedliche Werkzeuge zurückgreifen. Alle gängigen Bibliotheken aus \autoref{tab:dilist} liefern bereits Werkzeuge für die Konkatenierung mit.
	
	\item[AngularJS] Alle AngularJS-Templates, welche sich im Normalfall in einer jeweils eigenen HTML-Datei befinden, lassen sich in eine eigene JavaScript-Datei umzuwandeln. Als Werkzeug wird hier Grunt in Verbindung mit dem \ac{npm}-Paket namens \quotes{grunt-angular-template} \cite{Clemmons2016} verwendet. Dies nimmt die AngularJS-Templates und wandelt diese in eine JavaScript-Datei um, welche den Template-Cache \cite{Google2016c} von AngularJS nutzt.
	
	\item[AngularJS2] Im Gegensatz zum \ac{jit}-Ansatz, bei dem Skripte und App-Ressourcen erst dann geladen werden, wenn sie benötigt werden, wird beim \ac{aot}-Ansatz alles, was später benötigt werden könnte, zu Anfang geladen.  In der Dokumentation von AngularJS2 wird der \ac{aot}-Ansatz erklärt \cite{Google2016a} und wie dieser angewandt wird.
	
	\item[Aurelia] Auch Aurelia bietet ein derartiges Tool an. Dieses nutzt Gulp und ist über \ac{npm} mit dem Namen \quotes{aurelia-bundler} \cite{Aurelia2016} zu finden.
	
	\item[React] Um die Laufzeit der Webanwendung zu optimieren, kann man diese zuvor in eine normale JavaScript-Anweisung umwandeln. Hierfür empfiehlt sich der Einsatz von Babel \cite{Babel2016}. Das Werkzeug wird über \ac{npm} installiert, kann TypeScript und neueres ECMAScript in ECMAScript 3 und 5 umwandeln, und beherrscht auch \ac{jsx}. Die Umwandlung lässt sich manuell, oder auch wahlweise durch Gulp und Grunt anstoßen.
	
	
	
\end{description}
\subsection{Übertragen von Ressourcen}
Die Übertragung von Ressourcen in das \ac{aem} kann über mehrere Wege erfolgen.

\subsubsection{WebDAV}
Zunächst lässt sich über die \ac{webdav}-Schnittstelle mit einem entsprechenden Programm die Struktur des \ac{jcr} anzeigen und bearbeiten. Bei \acf{webdav}\index{WebDAV} handelt es sich um einen offenen Standard für die Bereitstellung von Daten über ein Netzwerk. Neben dem Löschen und Verschieben von Knoten ist hier auch das Hochladen von lokalen Dateien möglich.

\subsubsection{HTTP}
Weiterhin lassen sich Ressourcen auch über \ac{http} übertragen. Ein hilfreiches Programm hierfür ist \quotes{curl}\index{curl} \cite{Stenberg2017}. Über die Kommandozeile lassen sich so Dateien in das \ac{aem} übertragen. Folgender Befehl in \autoref{lst:curl} würde zum Beispiel eine Datei mit dem Namen \filefolder{flowerpot.png} in das \ac{jcr} unter \filefolder{/content/flowerpot.png} ablegen. Für die Authentifizierung wird ein gültiger Benutzername mit Passwort, hier beides \quotes{admin}, benötigt.

\begin{lstlisting}[style=jcr, caption=Eine Datei mit curl in das AEM übertragen, label=lst:curl]
curl -u admin:admin -T flowerpot.png http://localhost:4502/content/flowerpot.png
\end{lstlisting}


\subsubsection{CRX-Paket}
Sofern die \ac{webdav}-Schnittstelle auf der Zielplattform nicht geöffnet sein sollte, gibt es die Alternative, die App-Ressourcen als \ac{crx}-Paket bereit zu stellen. Hierfür werden zunächst wieder alle Inhalte in das \ac{jcr}, zum Beispiel über \ac{webdav}, geladen, dieses Mal jedoch in die der Entwicklungsumgebung. Anschließend kann über die Autoren-Bedienoberfläche, siehe \autoref{sec:autor_ui}, ein \ac{crx}-Paket erstellt werden. Dieses ist ein Zip-Archiv mit Metainformationen, wie Name und Version des \ac{crx}-Pakets.\\
Das Erstellen eines \ac{crx}-Paketes erfolgt über eine entsprechende Seite des CRXDE Lite. Hier werden besagte Metainformationen gesetzt und alle benötigten Knoten für das \ac{crx}-Paket\index{CRX-Paket} angegeben. Nun erfolgt der Export des \ac{crx}-Pakets. Ist dies geschehen, wird das \ac{crx}-Paket in die Zielplattform importiert. Der Inhalt des \ac{crx}-Paketes wird nun in das \ac{jcr} geschrieben. Bereits bestehende Knoten werden bei diesem Vorgang ggf. überschrieben.
Alternativ lässt sich das \ac{crx}-Paket auch mit Maven erstellen. Durch Maven\index{Maven} lassen sich standardisierte Anwendungen erstellen und verwalten. Eine entsprechende Anleitung zum Erstellen eines Maven-Projektes und wie sich hiermit ein \ac{crx}-Paket erschaffen lässt stellt Adobe zur Verfügung \cite{Incorporated2017};
\section{Suchmaschinen}
\label{sec:suchmaschinen}
Falls die Webanwendung als \ac{spa} mit Hinzunahme der in \autoref{sec:pushState} beschriebenen pushState-Funktion entwickelt wurde, so ist eine entsprechende Konfiguration des \ac{aem}-Servers vonnöten, damit die Webseiten der Webanwendung korrekt aufgerufen werden.\\
Als Beispiel soll eine AngularJS-Anwendung dienen, welche unter \filefolder{/content/angular/} liegt. Sie soll über \pseudourl{http://localhost:4502/spa/} zu erreichen sein. Somit sind alle App-Ressourcen der Webanwendung auch unter besagter \ac{url} erreichbar. \\
Es wird angenommen, dass alle \ac{url}s, die einen Punkt besitzen, einer App-Ressource wie Skripte, Bilder etc. entsprechen. Ansonsten handelt es sich um eine Webseite innerhalb der \ac{spa}. Da eine \ac{spa} aus einer HTML-Seite, z. B. index.html, besteht, sollen alle Webseiten-Anfragen auf diese verweisen. Alle weiteren App-Ressourcen erhalten eine relative Umleitung. \autoref{tab:resolution} zeigt beispielhaft einige vom Browser angeforderte Web-Ressourcen und unter welchen JCR-Knoten diese zu finden sind. \\

\begin{minipage}{\textwidth}
	\begin{longtable}{|p{0.40\textwidth}|c|c|p{0.25\textwidth}|}
		\hline  
		Angeforderte \ac{url}s & Ressourcetyp & Hat Punkt & Ziel JCR-Knoten\\  \hhline{|=|=|=|=|}
		
		\pseudourl{http://localhost:4502/spa/} & Webseite & Nein & \pseudourl{/content/angular/index.html} \\ 
		\hline
		 
		 \pseudourl{http://localhost:4502/spa/news/} & Webseite & Nein & \pseudourl{/content/angular/index.html} \\ 
		 \hline
		 
		 \pseudourl{http://localhost:4502/spa/about/} & Webseite & Nein & \pseudourl{/content/angular/index.html} \\ 
		 \hline
		 
		 \pseudourl{http://localhost:4502/spa/img/logo.png} & Bild & Ja & \pseudourl{/content/angular/img/logo.png} \\ 
		 \hline
		 
		 \pseudourl{http://localhost:4502/ws/news/} & Webseite & Nein & Keine Weiterleitung, da Pfad nicht mit \filefolder{/spa} beginnt.\\ 
		 \hline
		 
		 \pseudourl{http://localhost:4502/spa/ressources/style.css} & Stylesheet & Ja & \pseudourl{/content/angular/ressources/style.css} \\ 
		 \hline
		 
		\caption{Beispiel für die interne Umleitung}\label{tab:resolution}
	\end{longtable}
\end{minipage}
Hier kommt der Apache Sling Resource Resolver zum Einsatz \cite{Foundation2016}. Über das \ac{jcr} werden Regeln für die Umleitung der angeforderten \ac{url}s gesetzt. In der Standardkonfiguration befinden sich diese unter \filefolder{/http/map/}. Jede Regel setzt sich aus der hierarchischen Struktur und verschiedenen Eigenschaften zusammen. Um die anforderten \ac{url}s aus \autoref{tab:resolution} entsprechend umzuleiten, empfiehlt sich die Struktur wie in \autoref{lst:resolver}.

\begin{lstlisting}[style=jcr,caption=Konfigurationsbeispiel für den Apache Sling Resource Resolver, label=lst:resolver]
/etc/map/http/
+ localhost.4502/
  - jcr:primaryType (Name) = sling:Mapping
  + spa/
    - jcr:primaryType (Name) = sling:Mapping
    - sling:internalRedirect (String) = /content/angular
    + path/
      - jcr:primaryType (Name) = sling:Mapping
      - sling:match (String) = ([^.]*)
      - sling:internalRedirect (String) = /content/angular/index.html
\end{lstlisting}

Die Regeln basieren auf regulären Ausdrücken. Bei regulären Ausdrücken\index{Reguläre Ausdrücke} handelt es sich um eine Notation zur Beschreibung von Textmustern. Durch die Hinzunahme einer geeigneten Programmiersprache ist es möglich, Texte nach dem Textmuster zu durchsuchen und diese so zu erweitern, zu reduzieren und zu manipulieren \cite[S. 1 f.]{Friedl2009}. \\
Der reguläre Ausdruck wird aus den Namen der Knoten bzw. deren Werten der Eigenschaft sling:match, falls ein Knoten diese besitzt, zusammengesetzt. Der Wert der Eigenschaft sling:internalRedirct steht für die Ziel-URL im \ac{jcr}. \\
Aus der gegebenen Konfiguration resultieren zwei Regeln. Der Knoten \filefolder{/etc/map/http/localhost.4502/spa/path/} beschreibt, dass alle \ac{url}s, die mit \pseudourl{http://localhost:4502/spa/} beginnen, aber keinen Punkt in der \ac{url} haben, auf \filefolder{/content/angular/index.html} verweisen. Ansonsten gilt die erstellte Regel unter Knoten \filefolder{/etc/map/http/localhost.4502/spa/}.

Wie die Auflösung verschiedener \ac{url}s geschehen würde, ist der \autoref{img:resolve} zu entnehmen.

\begin{figure}[H]
	\begin{center}
		\includegraphics[width=1\textwidth]{resolver.png}
		\caption{Möglichkeit zum Auflösen von Anfragen unter AEM}
		\label{img:resolve}
	\end{center}
\end{figure}

Sollten die Ressourcen über einen Proxy geleitet werden, so ließen sich diese Regeln auch mit dem Modulen \quotes{mod\_rewrite} abbilden, sofern es sich bei dem Proxy um einen Apache HTTP-Server handelt.

\section{Bekannte Konflikte zwischen Skripten}
\label{sec:konflikte}
Während der Integration sind keinerlei Problematiken zwischen den Skripten von \ac{aem} und Skripten seitens der Webanwendung bzw. eines Webframeworks aufgetreten. \\
Auch beim Einsatz von unterschiedlichen Versionen des gleichen Webframeworks wurden keinerlei Probleme entdeckt.

\sectiontbd{Bewertung der Lösungen}

\missingall


\begin{longtable}{| c | c|} 
	\hline 
	\thead{Erfüllte Anforderung} & \thead{Gelöst in Kapitel} \\ 
	
	\hline 
	A1 &  \ref{sec:integrationen} \\ 
	\hline
	A2 &  \ref{sec:integrationen} \\  
	\hline
	A3 &  - \\  
	\hline 
	A4 &  \ref{sec:integrationen} \\   
	\hline
	A5 &  \ref{sec:integrationen} \\   
	\hline
	A5.1 &  \ref{sec:direkte-integration-in-das-aem}\\  
	\hline 
	A5.2 &  \ref{sec:ansatz-java} \\  
	\hline 
	A5.3 &  \ref{sec:ansatz-javascript}\\  
	\hline 
	A6 &  \ref{sec:aufbereiten} \\   
	\hline
	A7 &  \ref{sec:suchmaschinen} \\   
	\hline
	A8 &  \ref{sec:proxy}\\ 
	
	\hline 
	\caption{Anforderungserfüllung}\label{tab:erfuellung}
\end{longtable}


